%==============================================================================
%  README final / Rapport de projet  --  ST2MPG (Patrouille Forestiere Responsable)
%  Prototype : "Nettoyage de la Foret"
%  Compilation :  pdflatex RAPPORT.tex   (deux passes pour la table des matieres)
%==============================================================================
\documentclass[11pt,a4paper]{article}

\usepackage[utf8]{inputenc}
\usepackage[T1]{fontenc}
\usepackage[french]{babel}
\usepackage{lmodern}
\usepackage[margin=2.4cm]{geometry}
\usepackage{titlesec}
\usepackage{enumitem}
\usepackage{booktabs}
\usepackage{tabularx}
\usepackage{xcolor}
\usepackage{hyperref}

\definecolor{forest}{RGB}{28,90,170}
\hypersetup{
  colorlinks=true,
  linkcolor=forest,
  urlcolor=forest,
  citecolor=forest,
  pdftitle={README - ST2MPG Patrouille Forestiere Responsable}
}

\titleformat{\section}{\Large\bfseries\color{forest}}{\thesection}{1em}{}
\titleformat{\subsection}{\large\bfseries}{\thesubsection}{1em}{}
\setlist{itemsep=2pt, topsep=3pt}

% Champ à compléter (surligné en rouge) — remplacez le texte entre accolades.
\newcommand{\todo}[1]{\textcolor{red}{\textbf{[\,#1\,]}}}

\newcommand{\keycap}[1]{\fbox{\small\texttt{#1}}}

%==============================================================================
\begin{document}

%------------------------------------------------------------------ Page titre
\begin{titlepage}
\centering
\vspace*{2.2cm}
{\large ST2MPG --- Projet de Modélisation Géométrique\par}
{\large Fil rouge : \emph{Patrouille Forestière Responsable}\par}
\vspace{1.6cm}
{\Huge\bfseries\color{forest} Nettoyage de la Forêt\par}
\vspace{0.5cm}
{\Large Prototype interactif PC (clavier-souris) --- Unity\par}
\vspace{2cm}

{\large\bfseries Membres du groupe (trinôme)\par}
\vspace{0.4cm}
{\large
Junghoon Kim (LaelKim)\\[2pt]
\todo{Nom du membre 2}\\[2pt]
\todo{Nom du membre 3}\par}

\vfill
{\large Groupe \todo{n° de groupe} \quad---\quad \today\par}
\end{titlepage}

\tableofcontents
\newpage

%------------------------------------------------------------------ Résumé
\section{Résumé du prototype}
\emph{Nettoyage de la Forêt} est le prototype réalisé dans le cadre du fil rouge
\emph{Patrouille Forestière Responsable} (ST2MPG). Il met en scène une situation
de \textbf{surveillance forestière} à visée environnementale : le joueur, jeune
garde forestier, effectue une patrouille dans une réserve naturelle afin
d'\textbf{identifier et de traiter trois anomalies environnementales} --- un
\textbf{arbre malade}, un \textbf{animal} abandonné (chien) et un \textbf{sac de
déchets sauvages}. Le joueur repère chaque anomalie à proximité, la ramasse
(signalement + collecte) puis la rapporte à la \textbf{poubelle} près du poste
forestier avant la fin du \textbf{compte à rebours de 4 minutes}.

Un \textbf{garde forestier (PNJ)} patrouille entre les anomalies via le NavMesh et
salue le joueur pour l'orienter tant qu'une anomalie n'a pas été traitée. Le monde
est animé par un \textbf{cycle jour / nuit} dynamique, et l'expérience est encadrée
par un menu d'accueil (consignes de mission), un menu pause et un écran de bilan
final (mission réussie ou temps écoulé). Conformément au positionnement du module,
la priorité a été mise sur une \textbf{scène 3D structurée, cohérente et animée}
intégrant des assets modélisés sous Blender.

%------------------------------------------------------------------ Commandes
\section{Commandes clavier--souris}
\begin{center}
\begin{tabularx}{\linewidth}{@{}lX@{}}
\toprule
\textbf{Commande} & \textbf{Action} \\
\midrule
\keycap{Z}\,\keycap{Q}\,\keycap{S}\,\keycap{D} / Flèches & Déplacer le personnage \\
Souris & Orienter la caméra (troisième personne) \\
\keycap{Espace} & Sauter \\
\keycap{Maj gauche} & Courir / sprint \\
\keycap{E} & Ramasser (signaler) une anomalie proche / la déposer dans la poubelle \\
\keycap{Échap} & Ouvrir / fermer le menu pause \\
Clic gauche & Valider les boutons des menus (Commencer, Reprendre, Recommencer, Quitter) \\
\bottomrule
\end{tabularx}
\end{center}

\noindent\textit{Remarque : les déplacements et la caméra reposent sur le
contrôleur Invector (Third Person Controller LITE) ; l'interaction
(ramasser / déposer / signaler) et l'ensemble des menus sont des développements
propres au groupe.}

%------------------------------------------------------------------ Fonctionnalités
\section{Liste des fonctionnalités}
\begin{itemize}
  \item \textbf{Scène forestière jouable} : terrain naturel, végétation, poste
        forestier (maison en bois), zone de départ et PNJ animé.
  \item \textbf{Contrôleur troisième personne} : marche, course, saut et caméra
        libre à la souris pour explorer la zone.
  \item \textbf{Trois anomalies environnementales} identifiables visuellement
        (arbre malade, animal, déchets) à retrouver et à traiter.
  \item \textbf{Détection et signalement de proximité} : détection de l'anomalie
        la plus proche, message contextuel \og [E] Ramasser \fg{}, puis dépôt à la
        poubelle qui valide l'anomalie et met le score à jour.
  \item \textbf{Système de ramassage / transport} : animation de ramassage,
        attache de l'objet dans la main droite du personnage, transport, puis dépôt.
  \item \textbf{Compte à rebours} de 4 minutes visible dans l'interface ;
        l'affichage passe au rouge sous 30 secondes.
  \item \textbf{Conditions de fin claires} : victoire (3/3 déposées) ou défaite
        (temps écoulé), avec écran-bilan (temps utilisé, nombre d'anomalies).
  \item \textbf{Garde forestier (PNJ animé)} : patrouille automatique entre les
        anomalies via le NavMesh, se tourne vers le joueur et le salue tant que
        l'anomalie visée n'a pas été traitée (feedback / réaction à l'approche).
  \item \textbf{Cycle jour / nuit dynamique} : rotation du soleil, variation
        d'intensité, de couleur (blanc chaud, orange à l'horizon, bleu nuit),
        d'ambiance globale et d'exposition de la skybox.
  \item \textbf{Halo lumineux de nuit} : point de lumière pulsé autour du garde
        pour le repérer dans l'obscurité.
  \item \textbf{Interface utilisateur (HUD)} : barre supérieure avec minuteur,
        compteur \og 0 / 3 \fg{}, pastilles de progression et messages contextuels.
  \item \textbf{Menus} : écran d'accueil (rôle et objectifs), menu pause (gel du
        temps et du joueur) et écran de fin, avec gestion automatique du curseur.
\end{itemize}

%------------------------------------------------------------------ Assets modélisés
\section{Liste des assets modélisés}
Assets modélisés / texturés par le groupe sous Blender puis exportés en FBX et
intégrés dans Unity (échelle contrôlée, matériaux lisibles, nommage clair) :
\begin{itemize}
  \item \textbf{Arbre malade} (\texttt{ArbreMalade.fbx}) --- élément statique
        personnalisé servant d'anomalie.
  \item \textbf{Poubelle} (\texttt{trashcan.fbx}) --- objet environnemental, point
        de dépôt / validation des anomalies.
  \item \textbf{Maison en bois / poste forestier} (\texttt{Wood/Wood\_house.fbx})
        avec son jeu de textures PBR (base color, normal, roughness, metallic,
        displacement).
  \item \textbf{Garde forestier} --- texturage du personnage
        (\texttt{Garde\_Material} : diffuse camo, normal, roughness).
\end{itemize}

\noindent\textbf{Assets importés / externes} (non modélisés, utilisés tels quels
ou adaptés) : contrôleur \emph{Invector Third Person Controller LITE},
\emph{Supercyan Free Forest Sample} (décor forestier et personnage jouable),
modèle du chien \emph{Shiba} (animal-anomalie) et sac de déchets (\texttt{garbage}).

\noindent\textit{Couverture du cahier des charges de modélisation (§6) : deux objets
environnementaux (poubelle, sac de déchets), un bâtiment / poste forestier (maison
en bois), un élément statique personnalisé (arbre malade). L'animal-anomalie utilise
un modèle importé.}

%------------------------------------------------------------------ Animations
\section{Liste des animations utilisées}
\begin{itemize}
  \item \textbf{Locomotion du joueur} : idle, marche, course, saut (fournies par
        le contrôleur Invector).
  \item \textbf{Ramassage} (\texttt{Lifting.fbx}) : animation \og Pickup \fg{}
        jouée lorsque le joueur ramasse / signale une anomalie.
  \item \textbf{Marche du garde} (\texttt{Walking.fbx}) : déplacement du garde
        forestier (PNJ) pendant la patrouille.
  \item \textbf{Salut du garde} (\texttt{Waving.fbx}) : geste de la main déclenché
        lorsqu'il atteint une anomalie et fait face au joueur.
\end{itemize}
Ces clips sont pilotés par l'\emph{Animator Controller}
\texttt{GardeForestier\_Controller} (paramètres \texttt{Speed}, \texttt{isWaving})
pour le PNJ, et par le déclencheur \texttt{Pickup} pour le joueur.

%------------------------------------------------------------------ Répartition
\section{Répartition du travail}
Répartition selon les trois rôles du trinôme définis par le sujet (§9).
\begin{center}
\begin{tabularx}{\linewidth}{@{}p{3.2cm}p{3.4cm}X@{}}
\toprule
\textbf{Membre} & \textbf{Rôle} & \textbf{Contributions principales} \\
\midrule
\todo{Membre} & Rôle 1 --- Modélisation \& assets &
Modélisation Blender (arbre malade, poubelle, maison/poste forestier),
texturage du garde, exports FBX, intégration et échelle dans Unity. \\
\addlinespace
\todo{Membre} & Rôle 2 --- Personnage, rigging \& animation &
Intégration du garde forestier (PNJ), Animator Controller, clips Walking / Waving
et animation de ramassage, transitions et cohérence avec le scénario. \\
\addlinespace
Junghoon Kim (LaelKim) & Rôle 3 --- Intégration Unity, interaction \& UI &
Scène et organisation du projet, déplacement PC, détection / signalement des
anomalies, score, timer, cycle jour/nuit, HUD et menus, build et tests. \\
\bottomrule
\end{tabularx}
\end{center}
\noindent\textit{Ajustez l'attribution des rôles et complétez les noms selon les
contributions réelles (chaque membre doit avoir participé aux tests, aux commits
Git et à la conception du scénario).}

%------------------------------------------------------------------ Liens
\section{Liens}
\begin{itemize}
  \item \textbf{Dépôt Git} : \todo{URL du dépôt}
  \item \textbf{Vidéo de démonstration} : \todo{URL (YouTube non répertorié / Drive / OneDrive)}
  \item \textbf{Exécutable Windows} : \todo{URL de l'archive Build\_...\_GroupeX.zip}
\end{itemize}

%------------------------------------------------------------------ Bugs connus
\section{Bugs connus}
\begin{itemize}
  \item Le garde forestier dépend de la présence d'un \textbf{NavMesh} et
        d'objets nommés exactement \texttt{ArbreMalade}, \texttt{Shiba},
        \texttt{garbage} ; un renommage ou un NavMesh non généré interrompt sa
        patrouille.
  \item La détection de la poubelle et des anomalies se fait par \textbf{nom}
        d'objet ; toute duplication ou renommage peut empêcher le ramassage ou le
        dépôt.
  \item L'ajustement de l'échelle de l'objet tenu en main peut varier selon le
        modèle (pivot / unités du FBX), entraînant parfois un léger décalage à
        l'attache.
  \item Absence de \textbf{feedback sonore} (le feedback est uniquement visuel /
        textuel pour l'instant).
  \item \todo{Ajoutez ici tout autre bug constaté lors de vos tests}
\end{itemize}

%------------------------------------------------------------------ Améliorations
\section{Pistes d'amélioration}
\begin{itemize}
  \item Remplacer la recherche d'objets par nom par un système de \emph{tags} ou
        de composants dédiés, plus robuste.
  \item Ajouter du \textbf{son} : ambiance forêt jour/nuit, effets de ramassage /
        dépôt, message radio (animation d'utilisation de la radio par le garde,
        prévue par le sujet).
  \item Ajouter le \textbf{retour au poste forestier} comme condition de fin
        alternative, en plus de la collecte complète.
  \item Augmenter le nombre et la variété des anomalies, avec placement aléatoire
        pour la rejouabilité.
  \item Enrichir l'IA du garde (dialogue, pointage vers l'anomalie la plus proche).
  \item Ajouter un mini-radar / boussole indiquant les anomalies restantes.
  \item Options de configuration (sensibilité souris, volume, difficulté / durée).
  \item \textbf{Évolution vers la XR} l'année suivante : adapter les interactions
        (ramassage, menus) et les déplacements aux contrôleurs de casque.
\end{itemize}

\end{document}
